\section{Representations and subgroup projections}\label{sec:conventions}
Throughout, logs used in estimates are natural logs. We write \(G_L\) for the symmetric group of degree \(L\); spaces can equally well be indexed by sets of size \(L\). Denote the representation belonging to a partition \(\lambda\) by \(V_\lambda\), and its dimension by \(D_\lambda\). We also write \(V_\mu,D_\mu\) for an irreducible of a product group and its dimension. We take the group of degree 0 to be trivial, and allow empty partitions where appropriate. All representations are taken unitary (over \(\mathbb C\)). The trace and Schatten norms on a single matrix space use unnormalized trace unless otherwise indicated. We use \(\|M\|_p=(\operatorname{Tr}|M|^p)^{1/p}\), including operator norm at \(p=\infty\).

We use the following characteristic-zero facts and normalization conventions.
\begin{itemize}
\item We use tableau dimensions, the branching rule, and the hook-length formula for symmetric groups \cite[Theorems~2.6.5, 2.8.3, and~3.10.2]{Sagan2001}. Thus restriction to a symmetric group of smaller degree has shapes given by box deletions. Transposition of a partition corresponds to tensoring with the sign \cite[Theorem~6.7]{James1978}, and induction from two factors of a Young subgroup uses the Littlewood--Richardson coefficients \cite[Theorem~16.4]{James1978}.
\item We will also use ordinary Schur--Weyl duality on a complex tensor power, with the unitary group action on the factors and permutations of slots \cite[Theorem~4.57 and Corollary~4.59]{etingof}. For a local dimension \(h\), the factors in the decomposition for tensor degree \(l\) are the polynomial unitary-group representation with highest weight \(\gamma\) and \(V_\gamma\), for \(\gamma\vdash l\) of length at most \(h\). We use Schur characters in their semistandard-tableau form \cite[Definition~7.10.1]{Stanley1999}. In particular a positive operator acting identically in the slots has, on type \(\gamma\), largest eigenvalue given by the monomial of exponents \(\gamma\) on its eigenvalues sorted in decreasing order. This follows also by column strictness in the tableaux and the presence of the highest weight. The statement for zero eigenvalues can be taken by a limit.
\item Given an irreducible \(\mu\) of a subgroup \(R\) and a matrix \(M\) on \(V_\mu\), we can use a group algebra element of \(R\) with matrix \(M\) there, and zero matrix on its other types, viewed also in any larger group. Its coefficient at \(g\in R\) is
\[
 \frac{D_\mu}{|R|}\operatorname{Tr}\big(M \rho_\mu(g^{-1})\big),
\]
where \(\rho_\mu\) denotes the representation. We refer to the images of this element as matrix actions supported on \(\mu\). For \(M\) the orthogonal projector onto a unit vector \(e\), denote such an image by \(P_{\mu,e}\), specifying or taking as understood the ambient representation. On restriction to \(R\), this is an orthogonal projection using one direction of the carrier and acting identically on the corresponding multiplicity space. A full isotypic projection instead uses \(I_{V_\mu}\). We use the standard orthogonality and Plancherel formulas. In particular the sum of squared absolute values of the above coefficients is \(D_\mu\operatorname{Tr}(M^*M)/|R|\).
\end{itemize}

Write \(K_d(\lambda)\) for the Fourier matrix (expectation of the representing matrix) on \(V_\lambda\), \(\lambda\vdash N\), \(N=2^d\). The length of any partition we need to consider with \(K_d(\lambda)\ne 0\) is at most \(N/2\): one matching of switches already projects to invariants of a Young subgroup with \(N/2\) parts of size 2. The corresponding permutation module only has partitions of length at most \(N/2\) (for example, it occurs in an ordinary slot permutation tensor power with local dimension \(N/2\)).


\part{Weighted contraction}\label{part:weighted}

\section{Dimensions away from a long first row}
\label{ten:weighted:part:53}
Write $k=N-\lambda_1$ for the number of boxes below the first row. The following dimension estimate will be used in all three moment arguments.
\begin{lemma}\label{lem:surviving-dimension}
There is an absolute constant $c>0$ such that, for every $N\ge2$ and every partition $\lambda\vdash N$, $\lambda\ne(N)$, with at most $N/2$ rows,
\[
 \log D_\lambda\ge c(N-\lambda_1).
\]
\end{lemma}
\begin{proof}
First take $N$ sufficiently large. If both the width and the height are less than $N/10$, every hook length is at most $N/5$, and the hook formula gives $\log D_\lambda\ge cN\ge ck$.
Otherwise work in $\lambda$ or its transpose so the first row has size $u\ge N/10$. Let $v=N-u$ be the number of boxes outside that row. In the original orientation $v=k$; in the transposed orientation the height assumption gives $v\ge N/2$. Retain that row and a subdiagram of size $l=\min(v,\lfloor u/2\rfloor)$ below it. In either orientation $l\ge c_1k$ for an absolute $c_1>0$ and large $N$. By branching its dimension is a lower bound for $D_\lambda$. The hook formula bounds this dimension below by
\[
 \binom{u+l}{l}\exp\!\left(-\frac{l}{u-l+1}\right).
\]
Indeed the lower diagram has dimension at least one, its width is at most $l$, and its column heights sum to $l$, so the logarithm of the extra top-row hook factor is at most $l/(u-l+1)$. Since $l\le u/2$, the displayed lower bound has logarithm at least $(\log3-1)l\ge c_2k$. This proves the claim for large $N$. The remaining sizes are finite; decreasing $c$ covers them because the only one-dimensional irreducibles are the trivial and sign representations, both excluded by the hypotheses.
\end{proof}
In particular, for fixed small \(\epsilon>0\), if \(k>N^{1-\epsilon}\) and \(\lambda\) has length at most \(N/2\), then
\begin{equation}
 \log D_\lambda \ \ge\ c_0 N^{1-\epsilon}.                                                  \label{ten:weighted:eq:3}
\end{equation}


\section{Weights from hook trimming}
\label{ten:weighted:part:60}
The weights balance the dimension of a subdiagram against the number of boxes removed. Their parent--child difference will pay for the entropy of restoring deleted cells, while retaining a favorable dimension term. For $L\ge2$, fix
\[
 b_0=\tfrac14,\quad\delta=\tfrac1{16},\quad
 b_1=b_0\delta/4,\quad A=10,\quad
 \eta=b_1/(2A),\quad\epsilon=\eta/8,
\]
and define
\[
 F_L(r)=r\min\{\log(L/r),\eta\log L\}\quad(1\le r\le L),
 \qquad F_L(0)=0.
\]
For a Young diagram $\lambda\vdash L$, a subdiagram $a\subseteq\lambda$ is itself a Young diagram. Put
\begin{equation}\label{ten:weighted:eq:4}
 \begin{split}
 J(\lambda)&=\min_{a\subseteq\lambda}
 \{b_0\log D_a+b_1r\log L-AF_L(r)\},\qquad r=L-|a|,\\
 W_\lambda&=D_\lambda e^{-J(\lambda)}.
 \end{split}
\end{equation}
Empty diagrams have dimension one. For sizes zero and one put $J=0$ and $W=1$. Since $AF_L(r)\le(b_1/2)r\log L$, the expression minimized is nonnegative; taking $a=\lambda$ gives
$0\le J(\lambda)\le\frac14\log D_\lambda$.

On a tuple of partitions for a product of same-size groups, we sum the $J$ values and multiply the $W$ values. We next show that the minimizing subdiagram can be chosen in a hook at a controlled cost.

We will need that at \(L=N\), at additive cost at most \(O(h^2)\) in the minimum of \eqref{ten:weighted:eq:4}, we can take \(a\) in the \(h,h\)-hook, i.e. with no box of indices \(i>h,\ j>h\), for \(h=\lfloor N^\delta\rfloor\) and \(N\) sufficiently large. To show it, take a minimizing partition \(a_1\) and trim off the boxes with both indices \(>h\) leaving \(a\). Write \(x\) for the number trimmed and \(M=|a_1|\), supposing \(x>0\). The hook length formula gives
\begin{equation}
 \log D_{a_1}-\log D_a \ \ge\ x(\log h-C)-C h^2.                                            \label{ten:weighted:eq:5}
\end{equation}
Here the log factorial gain for \(a_1\) is at least \(x(\log M-1)\), and hooks in the trimmed boxes are at most \(2M/h\). Hooks within \(a\) changing length due to horizontal extensions are bounded in their log ratio sum by
\[
 \sum_{i=1}^{p}\sum_{j=1}^{h}\log\left(1+\frac{x_i}{h-j+1+p-i}\right),
\]
where \(p\) is the number of rows of the trimmed portion and \(x_i\) their lengths within that portion, decreasing weakly. If \(p\ge h\), linearize by \(\log(1+v)\le v\), and use that the sum of reciprocals over \(j\) increases with \(i\); its average over \(i\) is at most 2 (bound the double sum of reciprocals by extending to \(p\) columns and summing on diagonals). Thus pairing with decreasing \(x_i\) gives a bound \(2x\). If \(p<h\), the sum of the row contributions is at most
\[
 \sum_{i=1}^p\log\binom{h+x_i}{x_i}
 \le (ph+x)\log2\le(h^2+x)\log2.
\]
The transpose handles vertical extensions. This proves \eqref{ten:weighted:eq:5}. The change upwards in \(b_1 r\log N-A F_N(r)\) from increasing \(r\) by \(x\) is at most \((b_1\log N+A)x\), by the bound \(-1\) on how negative a slope of \(F_N\) can be. Since \(b_1<b_0\delta\), the asserted near-minimum property follows.

For such an \(a\), put \(t=N-|a|\). We will use
\begin{equation}
 b_0\log D_a + b_1 t\log N-A F_N(t)=J(\lambda)+O(h^2),\qquad
 t\log N\le C(\log D_\lambda+h^2),                                                        \label{ten:weighted:eq:6}
\end{equation}
where the second estimate uses the cap in \(F_N\). We allow the empty partition, with dimension 1.

We next prepare an estimate for transverse decompositions of a representation in this hook. This will later be used on a set of cells of a grid excluding \(t\) cells.


\section{Signed representations on occupied cells}
\label{ten:weighted:part:96}
\begin{theorem}[Overlap on an arbitrary occupied board]\label{thm:occupied}
Take a grid with \(m\) rows and \(n\) columns, \(mn=N\), and consider any subset of \(s=N-t\) cells, with the group \(G_s\) on those cells. Let \(a\vdash s\) lie in the \(h,h\)-hook, \(1\le h\le N\). Decompose \(V_a\) under the within-row permutation subgroup of \(G_s\), and also (separately) the within-column subgroup. For types \(\sigma,\theta\) of these two groups, let \(I_{\sigma,u}: V_\sigma\to V_a\), \(I_{\theta,v}: V_\theta\to V_a\) be isometric embeddings giving orthogonal copies and spanning the respective isotypic spaces. Then
\begin{equation}
 \sum_{u,v}\|I_{\sigma,u}^* I_{\theta,v}\|_\infty^2
 \ \le\
 \exp\{C[(m+n+1)h^2\log(N+2)+t]\}\ \min\left\{1,\frac{D_\sigma D_\theta}{D_a}\right\}.       \label{ten:weighted:eq:7}
\end{equation}
Here \(\sigma\) and \(\theta\) can each consist of a tuple of partitions, and the infinity norm is the operator norm. For \(s=0\) the empty representation has dimension one.
\end{theorem}
\begin{proof}
For \(s=0\) the assertion is immediate.

The signed realization uses ordinary Schur--Weyl duality \cite[Theorem~4.57 and Corollary~4.59]{etingof}, sign twist and induction \cite[Theorems~6.7 and~16.4]{James1978}. To prove the estimate we use a slot space
\(\mathcal L=(\mathbb C^h\oplus\mathbb C^h)^{\otimes s}\), calling the two blocks even and odd. Use the signed permutation action: on permuting a tensor of definite parities give the sign of reordering the subsequence of odd slots. This defines a group action (the signs compose), commuting with the global \(U(h)\times U(h)\) action for the two blocks. For clarity, the structure needed here follows from ordinary Schur-Weyl duality. Separate out the slots of even and odd parity. At fixed parity counts \(s_e,s_o\), the action of \(G_s\) moves the allocation of those counts to the slots, and at a fixed allocation the stabilizer action is unsigned permutation on the even slots and sign-twisted permutation on the odd slots. Consequently a \(U(h)\times U(h)\) sector of types \(\gamma,\xi\) (\(|\gamma|=s_e,|\xi|=s_o\), both lengths \(\le h\)) has multiplicity space giving, as a \(G_s\) representation,
\begin{equation}
 \operatorname{Ind}_{G_{s_e}\times G_{s_o}}^{G_s}\big(V_\gamma\otimes V_{\xi^{\rm tr}}\big).   \label{ten:weighted:eq:8}
\end{equation}
Here \(\xi^{\rm tr}\) denotes the transposed partition, and we use ``sector'' with its unitary-group carrier as well as the multiplicity space. There is such a sector containing \(V_a\): take \(\gamma\) from the first \(h\) rows and \(\xi^{\rm tr}\) the remaining rows (a diagram of width at most \(h\)). Indeed \(a/\gamma\) is a straight partition up to shift, so the Littlewood-Richardson coefficient in \eqref{ten:weighted:eq:8} is positive.

Here are size and entropy estimates used with these sectors, applying just as well at any number \(l\le s\) of slots. The dimensions of the unitary-group carriers and the number of sectors are \(\exp(O(h^2\log(N+2)))\), as upper bounds; for dimensions this follows for example by counting the possible counts of each symbol in each row of a semistandard tableau. The multiplicity of any \(\chi\vdash l\) for the signed permutation action on the \(l\) slots is also at most \(\exp(O(h^2\log(N+2)))\). Indeed, for a pair \(\gamma,\xi\) giving a nonzero induction coefficient at \(\chi\), use Littlewood-Richardson fillings of \(\chi/\xi^{\rm tr}\) with content \(\gamma\), and thus entries at most \(h\). Since \(\xi^{\rm tr}\) has width at most \(h\), every skew column beyond column \(h\) is a full column of \(\chi\). Column strictness therefore places these columns within the first \(h\) rows, where the fillings are determined by row counts of symbols. In the first \(h\) columns the skew shape has at most \(h^2\) boxes. These bounds along with \eqref{ten:weighted:eq:8} give the multiplicity estimate.

If \(|\gamma|+|\xi|=l\), write
\[
 H(\gamma,\xi)=\sum_i\gamma_i\log(l/\gamma_i)+\sum_j\xi_j\log(l/\xi_j)
\]
with zero contributions understood at count 0. For the same occurrence of \(\chi\), we have
\begin{equation}
 D_\chi\ \le\ \exp(H(\gamma,\xi))\ \le\ D_\chi\exp(C h^2\log(N+2)).                         \label{ten:weighted:eq:9}
\end{equation}
For the first bound use \(D_\chi\le\binom{l}{|\gamma|}D_\gamma D_\xi\), by \eqref{ten:weighted:eq:8} and invariance of dimension under transposition, and bound each factor by the multinomial of its row lengths. For the second bound, \(\chi\) is in the \(h,h\)-hook and contains both \(\gamma\) and \(\xi^{\rm tr}\). Since \(|\gamma|+|\xi^{\rm tr}|=|\chi|\), the number of boxes of \(\chi\) outside their union equals \(|\gamma\cap\xi^{\rm tr}|\), which is at most \(h^2\). In particular the total excess of its first \(h\) row lengths over the corresponding \(\gamma_i\) is \(O(h^2)\), and similarly for the first \(h\) column lengths and \(\xi_j\). In bounding its product of hooks, the part in its first \(h\) rows beyond column \(h\) has leg lengths bounded by \(h\), and the analogous statement for arms applies below its first \(h\) rows. This gives an upper bound on the hook product of
\[
 (N+2)^{C h^2}\prod_i\gamma_i!\prod_j\xi_j!
\]
(for example, for the horizontal part use the top row factorials, increasing their arguments by at most \(h\) in each row, and likewise for columns in the vertical part). Stirling's bound, or the elementary multinomial entropy bounds, now gives \eqref{ten:weighted:eq:9}.


\subsection{Two bounds for transverse overlap}
\label{ten:weighted:part:127}
We now prove the two bounds whose minimum occurs in Theorem~\ref{thm:occupied}. Counting copies gives the first; positive states and two-sided normalization give the dimension ratio.

We already get a bound \(\exp(O((m+n)h^2\log(N+2)))\) on the number of matrices \(I_{\sigma,u}^*I_{\theta,v}\), hence also on the sum in \eqref{ten:weighted:eq:7}, by using the multiplicity estimate row by row and column by column in \(\mathcal L\). When grouping slots for a subgroup, we can rearrange them with the signed reordering, after which contiguous block permutations give the tensor product actions. It remains to get the other bound in \eqref{ten:weighted:eq:7}. Let $P_\sigma,P_\theta$ be the isotypic projections on $V_a$. The all-copy identity
\[
 \sum_{u,v}\|I_{\sigma,u}^*I_{\theta,v}\|_2^2
 =\operatorname{Tr}_{V_a}(P_\sigma P_\theta)
\]
follows by expanding each projection as the sum over its orthogonal embeddings. We bound this sum of squared Hilbert--Schmidt norms, which also bounds the required sum of squared operator norms. Use a global unitary-group sector containing \(V_a\) as above, with projector \(P\) in \(\mathcal L\), and types denoted by \(\gamma,\xi\). If \(Q_\sigma,Q_\theta\) denote the subgroup isotypic projections on \(\mathcal L\), the trace we need to estimate is bounded by
\[
 \operatorname{Tr}_{\mathcal L}(Q_\sigma P Q_\theta P).
\]
Indeed \(P\) commutes with the permutations and each \(G_s\)-irreducible contribution to the trace on that sector is nonnegative.

We can majorize \(Q_\sigma\) in positive semidefinite order by a mixture of products \(S\) of states on individual slots, using the same state \(S_i\) in every slot of row \(i\), each state even (block diagonal). Here and below a state is positive semidefinite with trace 1. The scalar prefactor for this majorization can be taken to be \(D_\sigma\exp(Cm h^2\log(N+2))\), with the mixture a probability mixture. To see this on one row of size \(l>0\), bound the projection for its partition \(\chi\) by the sum of sector projections for compatible \((\gamma_0,\xi_0)\) at \(l\) slots. In each case twirl, by the unitary group \(U(h)\times U(h)\), the identical-slot product of a block diagonal state with diagonal entries \((\gamma_{0,i}/l,\xi_{0,j}/l)\). It acts after twirling as a scalar on this sector, at least \(\exp(-H(\gamma_0,\xi_0))\) divided by the dimension of the unitary-group carrier. Indeed before twirling it has the eigenvalue given by the highest weight and acts as identity on multiplicities in the sense of a matrix action on the carrier tensored with identity. This latter action of an identical even positive matrix in tensor power follows directly from ordinary Schur-Weyl on the two parities (and by limits for singular matrices). The twirl is positive on the other sectors. Now the bound on the prefactor follows from \eqref{ten:weighted:eq:9} and the counting and dimension bounds. A size-zero row requires nothing and can use any even state for its specification. The same majorization applies to \(Q_\theta\), with products \(T\) using \(T_j\) in column \(j\), and prefactor \(D_\theta\exp(C n h^2\log(N+2))\). Even when undoing a signed reordering of slots, such products are ordinary products since their factors preserve the local parity. Thus by two positive semidefinite majorizations (tracing with the other positive matrix conjugated by \(P\)), we can use bounds on
\(\operatorname{Tr}_{\mathcal L}(S P T P)\).

The normalization has two purposes: the global sector operator must contribute $\exp(-H(\gamma,\xi))$ to the squared norm, and the local overlaps must have mean at most one under the independent row and column marginals. Quarter-root filters meet both requirements. For fixed \(S,T\) so specified, let \(\bar S,\bar T\) be the average one-slot states over the \(s\) cells in consideration, and put \(X=\bar S^{1/2}\), \(Y=\bar T^{1/2}\). Write
\[
 S_i'=X^{-1/2} S_i X^{-1/2},\qquad T_j'=Y^{-1/2}T_jY^{-1/2}
\]
with inverses on supports. Let \(S',T'\) be the corresponding products, so the filtering can be undone using \(X^{1/2},Y^{1/2}\) on occupied slots. With
\[
 L=(X^{1/2})^{\otimes s} P (Y^{1/2})^{\otimes s}
\]
we have
\[
 \operatorname{Tr}_{\mathcal L}(S P T P)=\|(S')^{1/2} L (T')^{1/2}\|_2^2 .
\]
The matrix action \(L\) is supported on the \((\gamma,\xi)\) unitary-group sector with operator norm at most \(\exp(-H(\gamma,\xi)/2)\). Indeed \((X^{1/2})^{\otimes s}=(\bar S^{1/4})^{\otimes s}\) acts on this sector with norm at most \(\exp(-H(\gamma,\xi)/4)\), by taking largest weight monomials on the eigenvalues and using that the eigenvalues of \(\bar S\) sum to 1. The same argument applies with \(\bar T\). These matrix actions are identical on equivalent copies for the block unitary group. By matrix coefficient orthogonality for that compact group, \(L\) can therefore be represented by an integral of the global \(U^{\otimes s}\) for block unitaries \(U\), with coefficient function bounded by \(\|L\|_\infty\) times the squared carrier dimension. (This is the matrix Fourier formula for an action supported on one type.)

For each block unitary define \(z_{ij}=\operatorname{Tr}(S_i' U T_j' U^*)\). The squared Hilbert-Schmidt norm weighted by \(S',T'\), as above but for \(U^{\otimes s}\), is the product of \(z_{ij}\) over the cells under consideration. If \(p_i,q_j\) are the marginals of a uniformly chosen cell among these \(s\), the mean of \(z_{ij}\) under \(p_i q_j\) is at most 1. In fact it is \(\operatorname{Tr}(\bar S^{1/2} U\bar T^{1/2}U^*)\), at most 1 by Cauchy-Schwarz. Under the actual uniform law on the cells the mean of \(\log z_{ij}\) is thus bounded above by the relative entropy of that law with respect to \(p_i q_j\), by the elementary entropy variational inequality (the entropy/product duality is discussed in \cite[Theorem~2.1 and (2.4)]{CarlenCorderoErausquin2009}). Explicitly, if $\omega$ is the occupied-cell law and $r_{ij}=p_iq_j$, Jensen applied to $\sum\omega_{ij}\log(z_{ij}r_{ij}/\omega_{ij})$ gives at most $\log\sum r_{ij}z_{ij}\le0$. Zero values in the product need no bound. The relative entropy is at most \(\log(mn/s)\), by bounding the two marginal entropies by \(\log m,\log n\), so the product is at most \(\exp(s\log(N/s))\le\exp(t)\). Using the integral for \(L\) and the triangle inequality, we obtain
\[
 \operatorname{Tr}_{\mathcal L}(S P T P)
 \ \le\ \exp\{C h^2\log(N+2)+t\}\exp(-H(\gamma,\xi)).
\]
Since \(D_a\le\exp(H(\gamma,\xi))\), the mixture prefactors give the required bound with \(D_\sigma D_\theta/D_a\) in \eqref{ten:weighted:eq:7}, completing the proof.
\end{proof}


\section{The row-column Fourier estimate}
\label{ten:weighted:part:157}
Split a sweep of dimension \(d\) into two parts, each scanning consecutive coordinates, with block sizes \(d_1=\lfloor d/2\rfloor\) and \(d_2=d-d_1\). Index positions as a grid, so these parts use the within-row and the within-column subgroups, denoted \(R,C_{\rm col}\). Write \(n\) for the size of a row and \(m\) for that of a column, so \(mn=N\), \(m,n\) both within a constant factor of \(\sqrt N\). We take \(d\) large. Suppose henceforth in the estimate that
\begin{equation}
 k=N-\lambda_1>N^{1-\epsilon},\qquad \lambda^{\rm tr}_1\le N/2.                            \label{ten:weighted:eq:10}
\end{equation}

For irreducibles \(\mu,\nu\) of \(R,C_{\rm col}\), each a tuple, and unit carrier vectors \(e,f\), we claim the following bound on projections in \(V_\lambda\):
\begin{equation}
 W_\lambda\operatorname{Tr}_{V_\lambda}(P_{\mu,e}P_{\nu,f})
 \ \le\ W_\mu W_\nu\exp\{O((\log D_\lambda)/d)\}.                                          \label{ten:weighted:eq:11}
\end{equation}
All constants are independent of the types and of \(d\).

Use the hook \(a\) from \eqref{ten:weighted:eq:6} and choose a partition \(b\vdash t\) such that
\[
 \mathcal I=\operatorname{Ind}_{G_{N-t}\times G_t}^{G_N}(V_a\otimes V_b)
\]
contains \(V_\lambda\), possible by branching and decomposing the restriction. We bound the trace of \(X Y\) on \(\mathcal I\), for \(X=P_{\mu,e}\), \(Y=P_{\nu,f}\) acting here; this bounds the desired trace before multiplying by \(W_\lambda\), as the trace on each contributing \(G_N\)-type is nonnegative. Use the direct sum structure of \(\mathcal I\) with a fiber at each \(t\)-set \(Z\) of the grid. Its stabilizer acts on this fiber by \(V_a\) on the complement and \(V_b\) on \(Z\).

Let \(v_i\) be the counts of \(Z\) in rows and \(u_j\) in columns. The full row and column orbit sizes and overall size we denote by
\[
 L_R=\prod_i\binom n{v_i},\qquad L_C=\prod_j\binom m{u_j},\qquad L_0=\binom Nt.
\]
Let \(L_{RC}\) be the number of sets with the indicated two profiles. Use untruncated quantities
\[
 H_R=\sum_i v_i\log(n/v_i),\quad H_C=\sum_j u_j\log(m/u_j),\quad H_0=t\log(N/t),
 \quad E_0=H_R+H_C-H_0.
\]
Here again terms at zero count are zero. For nonempty layout classes \(E_0\ge 0\); for \(t>0\) it is \(t\) times the relative entropy of the uniform distribution on \(Z\) with respect to the product of its marginals. We have
\begin{equation}
 \log \frac{L_{RC}^2 L_0}{L_R L_C}\le E_0+O(t).                                            \label{ten:weighted:eq:12}
\end{equation}
Indeed \(\log L_R\ge H_R,\log L_C\ge H_C,\log L_0\le H_0+t\) by binomial bounds. For a uniform random set with the two profiles, subadditivity bounds \(\log L_{RC}\) by the sum of marginal indicator entropies. To bound that sum, choose a cell uniformly and call its indices \(I,J\) and its random indicator \(V\). As \(I,J\) are independent,
\[
 H(V\mid I,J)\le H(V\mid I)+H(V\mid J)-H(V)
\]
in ordinary entropy notation (by nonnegativity of conditional mutual information). Multiplying by \(N\), the terms on the right give an upper bound at most \((H_R+t)+(H_C+t)-H_0\), since the complementary term to \(v_i\log(n/v_i)\) in row binary entropy times row size is at most \(v_i\), and similarly for columns. Thus $\log L_{RC}\le E_0+2t$. Combining this with the three binomial bounds gives
\[
 2\log L_{RC}+\log L_0-\log L_R-\log L_C\le E_0+5t,
\]
which proves \eqref{ten:weighted:eq:12}.


\subsection{Positive blocks and stabilizer coefficients}
\label{ten:weighted:part:194}
The block trace for \(X Y\) only uses pairs \(Z,Z'\) with both profiles the same between them. In bounding the absolute sum of block contributions for one layout class, we can use \(L_{RC}^2\) times the maximum diagonal-block overlap \(\operatorname{Tr}(X_{ZZ}Y_{ZZ})\) in that class. In fact an off-diagonal block of a positive matrix has the form \(X_{ZZ}^{1/2} U X_{Z'Z'}^{1/2}\) for a contraction \(U\), and likewise for \(Y\); thus \(|\operatorname{Tr}(X_{ZZ'}Y_{Z'Z})|\) is bounded using Hilbert-Schmidt by the geometric mean of the corresponding diagonal-block overlap traces.

For fixed \(Z\), the row stabilizer \(R_Z\) in \(R\) acts separately on its sites outside and inside \(Z\). Write its types as \(\alpha=(\sigma,\tau)\) for the two parts respectively. The diagonal block \(X_{ZZ}\) has the following formula: on the fiber it is
\begin{equation}
 \frac{D_\mu}{L_R}\sum_{\alpha}\frac{1}{D_\alpha}\, \mathcal P_{\alpha,M_\alpha}.             \label{ten:weighted:eq:13}
\end{equation}
Here \(\mathcal P_{\alpha,M_\alpha}\) denotes a matrix action supported on \(\alpha\) for \(R_Z\), where the \(M_\alpha\) are positive matrices of total trace 1. More explicitly, write
\[
 V_\mu\!\downarrow_{R_Z}=\bigoplus_\alpha V_\alpha\otimes\mathcal M_{\mu,\alpha},\qquad
 e=\bigoplus_\alpha e_\alpha,\qquad
 M_\alpha=\operatorname{Tr}_{\mathcal M_{\mu,\alpha}}|e_\alpha\rangle\langle e_\alpha|.
\]
Then $M_\alpha\succeq0$ and $\sum_\alpha\operatorname{Tr}M_\alpha=\sum_\alpha\|e_\alpha\|^2=1$. Keeping just permutations in $R_Z$ in the matrix coefficient formula for $P_{\mu,e}$, and using $|R|=L_R|R_Z|$, gives \eqref{ten:weighted:eq:13} by orthogonality. The multiplicity spaces have been traced out, not counted. Similarly for \(Y_{ZZ}\) we use column stabilizer types \(\zeta=(\theta,\pi)\) with factor \(D_\nu/L_C\). Only types occurring in the corresponding restrictions and the fiber need be considered.

The positive states in \eqref{ten:weighted:eq:13} need not be product states on the surviving and deleted carriers. The following elementary estimate is the reason this causes no extra dimension factor.
\begin{lemma}[Arbitrary carrier vectors]\label{lem:entangled}
Let \(A_i:E^\prime\to E\) and \(B_j:F^\prime\to F\) be finite families of operators between finite-dimensional Hilbert spaces. Suppose
\[\sum_j|\langle x,B_jy\rangle|^2\le C\|x\|^2\|y\|^2.\]
Then for unit \(w\in E\otimes F\), \(w^\prime\in E^\prime\otimes F^\prime\),
\[\sum_{i,j}|\langle w,(A_i\otimes B_j)w^\prime\rangle|^2\le C\sum_i\|A_i\|_\infty^2.\]
\end{lemma}
\begin{proof}
For fixed \(i\), the functional \(B\mapsto\langle w,(A_i\otimes B)w^\prime\rangle\) has norm at most \(\|A_i\|_\infty\) on the operator-norm space. By trace-norm duality and singular-value decomposition it is a sum \(\sum_s c_s\langle x_s,By_s\rangle\), where \(x_s,y_s\) are unit vectors, \(c_s\ge0\), and \(\sum_s c_s\le\|A_i\|_\infty\). The triangle inequality in \(\ell^2(j)\) bounds the corresponding square root by \(\sqrt C\sum_s c_s\). Square and sum over \(i\).\end{proof}

For these fiber actions, take any pure unit states on the carriers \(V_\sigma\otimes V_\tau\) and \(V_\theta\otimes V_\pi\), denoted by \(w,w'\) as vectors. The corresponding overlap trace on the fiber is of the form
\begin{equation}
 \sum_{l,r} |\langle w,(A_l\otimes B_r) w'\rangle|^2.                                      \label{ten:weighted:eq:14}
\end{equation}
Here \(A_l\) range over the products \(I_{\sigma,u}^* I_{\theta,v}\) for the two decompositions of \(V_a\) as in \eqref{ten:weighted:eq:7}, and \(B_r\) analogously use \(V_b\), for types \(\tau,\pi\). This follows by writing the projection directions in every pair of copies in the restricted representations.

On \(Z\) itself, the within-row and within-column groups have trivial intersection. For unit vectors \(x,y\) of the two relevant carriers there, Plancherel gives
\begin{equation}
 \sum_r |\langle x,B_r y\rangle|^2
 \ \le\ \frac{D_\tau D_\pi}{D_b}\frac{t!}{\prod_i v_i!\prod_j u_j!}
 \ \le\ \frac{D_\tau D_\pi}{D_b}\exp(E_0+O(t)).                                             \label{ten:weighted:eq:15}
\end{equation}
Indeed the sum is a trace of a product of projections on \(V_b\), hence also the squared Hilbert-Schmidt norm of the product. The squared group-coefficient sum for the product of the two subgroup matrix elements used to give the projections (as group algebra elements) is the product of their squared coefficient sums, since the pairwise products of group elements use unique factorizations from those two subgroups. This gives the first bound using Plancherel in \(G_t\). The log factorial ratio is bounded as stated using \(E_0=t\log t-\sum_i v_i\log v_i-\sum_j u_j\log u_j\) and elementary factorial bounds.

Lemma~\ref{lem:entangled} bounds \eqref{ten:weighted:eq:14} by the right side of \eqref{ten:weighted:eq:15} times \(\sum_l\|A_l\|_\infty^2\), for arbitrary, possibly entangled, carrier vectors. We next keep the gain from the minimum in Theorem~\ref{thm:occupied}; it is needed when the placement entropy is paid.

Denote
\[
 S=\max\{0,\log(D_\sigma D_\theta/D_a)\}.
\]
Using \eqref{ten:weighted:eq:7}, \eqref{ten:weighted:eq:14}, and \eqref{ten:weighted:eq:15}, after dividing by \(D_\alpha D_\zeta\) from the diagonal-block formulas, the pure-state bound becomes
\begin{equation}
 \frac{1}{D_aD_b}
 \exp\{E_0-S+O(t+(m+n+1)h^2\log(N+2))\}.                                                  \label{ten:weighted:eq:16}
\end{equation}
The matrices within \eqref{ten:weighted:eq:13} and the column counterpart are handled by positive combinations of the pure states with total weights 1. Thus we can take upper bounds for the possible stabilizer types without a count-factor here. Multiplying by \(D_\mu D_\nu L_{RC}^2/(L_R L_C)\), and using \(D_\lambda\le L_0 D_a D_b\) and \eqref{ten:weighted:eq:12}, bounds the overlap for each layout class using \eqref{ten:weighted:eq:16}, by a bound of the form
\begin{equation}
 \frac{D_\mu D_\nu}{D_\lambda}
 \exp\{2E_0-S+O(t+(m+n+1)h^2\log(N+2))\},                                                 \label{ten:weighted:eq:17}
\end{equation}
taking the largest needed over contributing types.


\subsection{The budget pays for layout entropy}
\label{ten:weighted:part:234}
The restored-cell estimate still contains the cost \(2E_0-S\). We now compare it with the difference between the parent and child diagram weights.

To apply \eqref{ten:weighted:eq:4}, each \(\sigma_i\) for a complement row of size \(n-v_i\) is a subpartition of \(\mu_i\), due to restriction, and likewise \(\theta_j\subseteq\nu_j\) has size \(m-u_j\). Hence, denoting the sums of \(J\)'s by \(J(\mu),J(\nu)\), we have
\[
 J(\mu)+J(\nu)\ \le\
 b_0\log(D_\sigma D_\theta)+b_1 t\log N
 -A\left(\sum_i F_n(v_i)+\sum_j F_m(u_j)\right).
\]
Together with \eqref{ten:weighted:eq:6} this gives
\begin{equation}
 J(\lambda)-J(\mu)-J(\nu)
 \ \ge\ -b_0 S + A E_0 -O\big(h^2+(m n^{1-\eta}+n m^{1-\eta})\log(N+2)\big).               \label{ten:weighted:eq:18}
\end{equation}
Indeed the cap in the parent \(F_N\) term does not hurt this lower bound, and the amounts dropped from the row contributions to \(H_R\) come only from positive \(v_i\) less than \(n^{1-\eta}\), with a similar bound for \(H_C\).

Indeed, if $\mathcal E$ denotes the error on the right of \eqref{ten:weighted:eq:18}, then
\[
 2E_0-S-[J(\lambda)-J(\mu)-J(\nu)]
 \le (2-A)E_0-(1-b_0)S+\mathcal E\le\mathcal E.
\]
Here $A\ge2$, $b_0\le1$, and $E_0,S\ge0$. Thus the negative dimension term survives the off-diagonal count and is exactly strong enough when the ratio $D_\mu D_\nu/D_\lambda$ is replaced by $W_\mu W_\nu/W_\lambda$. All the log errors so far and the log count of possible profile choices, the latter \(O((m+n)\log(N+2))\), are bounded by \(O((\log D_\lambda)/d)\). Indeed \(t\) satisfies \eqref{ten:weighted:eq:6}, \(h\le N^{1/16}\), \(m,n\) are of order \(\sqrt N\), and we have \eqref{ten:weighted:eq:3} with \(\epsilon<\eta/2\). These estimates are for all sufficiently large \(d\) under \eqref{ten:weighted:eq:10}. This proves \eqref{ten:weighted:eq:11}.


\section{Weighted Schatten interpolation}
\label{ten:weighted:part:249}
The row part and the column part of the sweep, as group algebra elements, multiply to give \(K_d\) in the order specified by the scan (either product order will be fine for estimates). Within each row, the row part is an independent sweep in its smaller number of bits, and analogously for the columns. This holds because a group of consecutive bit updates in this sweep uses matchings confined to the indicated sets of positions, independently across the sets.

Still under \eqref{ten:weighted:eq:10}, let \(\mathcal B(M,M')\) on \(V_\lambda\) be the product of matrix actions, one for each subgroup, using arbitrary matrices \(M_\mu\) and \(M'_\nu\) on their respective irreducibles. Each subgroup action here is the sum over its irreducibles. By \eqref{ten:weighted:eq:11} we have
\begin{equation}
 W_\lambda \|\mathcal B(M,M')\|_2^2
 \ \le\ \exp\{O((\log D_\lambda)/d)\}
 \left(\sum_\mu W_\mu\|M_\mu\|_2^2\right)\left(\sum_\nu W_\nu\|M'_\nu\|_2^2\right).          \label{ten:weighted:eq:19}
\end{equation}
For a single pair of types this follows by writing the squared norm as a trace of a product of positive matrix actions (using the matrices times their adjoints or the reverse, as appropriate), then diagonalizing and using \eqref{ten:weighted:eq:11}. To sum in \eqref{ten:weighted:eq:19} use the triangle and Cauchy-Schwarz inequalities, absorbing log count factors into the given error. The log counts of the tuples of partitions for the two groups are \(O(N^{3/4}\log(N+2))\), using e.g. the partition count bound \(\exp(O(\sqrt n\log(n+2)))\) (split a partition according to parts of size greater than \(\sqrt n\) and the rest). So \eqref{ten:weighted:eq:3} suffices. Also
\[
 \|\mathcal B(M,M')\|_\infty\le \max_\mu\|M_\mu\|_\infty\ \max_\nu\|M'_\nu\|_\infty
\]
by unitary restriction. For \(2\le p<\infty\), interpolation therefore gives
\begin{equation}
 W_\lambda \|\mathcal B(M,M')\|_p^p
 \ \le\ \exp\{O((\log D_\lambda)/d)\}
 \left(\sum_\mu W_\mu\|M_\mu\|_p^p\right)\left(\sum_\nu W_\nu\|M'_\nu\|_p^p\right),          \label{ten:weighted:eq:20}
\end{equation}
with the coefficient in the log error bounded independently of \(p\).

For details, let $E=C(\log D_\lambda)/d$ be a common logarithmic error bound in \eqref{ten:weighted:eq:19}. The direct-sum Schatten norms use traces weighted blockwise by the given positive numbers and the output norm similarly. The bilinear map has norm at most $e^{E/2}$ at the $2$ endpoint and at most one at the $\infty$ endpoint. Normalize both inputs to norm 1 at \(p\), and test the output using the output weighted trace with a dual-norm-one matrix at \(p'=p/(p-1)\). In the strip with value taken at \(z=2/p\), extend the inputs analytically blockwise using their polar factors and exponent \(p z/2\) on their absolute values (zero singular directions may be kept zero). Their norms on the two boundaries, at \(\infty\) and 2 respectively, are then bounded by 1. Extend the matrix used in the trace test similarly (analytic for the matrix appearing in the trace), with exponent \(p'(1-z/2)\) for boundary norms at 1 and 2. At $\theta=2/p$, the three-lines inequality applied to the weighted trace with the bilinear product gives the factor $e^{(E/2)\theta}=e^{E/p}$. Taking the $p$-th power gives $e^E$ in \eqref{ten:weighted:eq:20}, with no growth of the logarithmic error coefficient in $p$.


\section{Sparse representations and completion}\label{sec:completion}
We have established the matrix inequality needed for the recursive range. A direct estimate handles representations whose first row contains almost all boxes. We prove that estimate here, then close the induction.
\subsection{A sparse forest estimate}
\label{ten:weighted:part:21}
Write \(k=N-\lambda_1\). For \(k\ge 1\) we have the bound
\begin{equation}
 \|K_d(\lambda)\|_\infty \ \le\ (C d)^k (k/N)^{k/8}                                                   \label{ten:weighted:eq:1}
\end{equation}
for an absolute \(C\). We prove it using the sweep kernel \(K\) on ordered injective \(k\)-tuples. Consider the subspace of functions \(f\) of the tuples with sum zero over any one component when all the others are fixed. This space is invariant for the position-permutation action. The module on the tuples contains \(V_\lambda\) by branching (invariants upon restriction to \(G_{N-k}\)). Since \(V_\lambda\) does not occur on \((k-1)\)-tuples, the isotypic space for \(\lambda\) on \(k\)-tuples belongs to the indicated zero-sum subspace. The kernel acts using the sweep in this representation or its adjoint convention.

In a transition \(x\mapsto y\) of tuples, a path for each component is specified just by its start and end: the successive coordinates at the end replace those of the start. For two components, let \(l\) be the last differing bit at the start, and \(r\) the first differing bit at the end. Give this pair weight \(w=0\) if \(r>l\), weight 2 if \(r=l\), and weight 1 otherwise. Then
\[
 K(x,y)=N^{-k}\prod_{\text{pairs}}w.
\]
Indeed, the paths have no conflict (two at the same position at the same stage boundary) exactly when \(r\le l\) for all pairs. When there is no conflict, a switch used by two paths instead of one is used by a pair with \(r=l\), at that bit, giving the factor 2. The same formula works for any subset of components.

Acting on the specified functions \(f\), we can replace the kernel by
\[
 D(x,y)=N^{-k}\sum_{B\subseteq [k]}(-1)^{k-|B|}\prod_{\text{pairs in }B}w,
\]
because all terms for proper subsets use zero marginal sums. This expression vanishes unless the graph of pairs with \(w\ne 1\) has no isolated vertex. Also, its absolute column sums are at most \(2^k\), since each subset term in absolute value uses the nonnegative sweep kernel on that subset (doubly stochastic), with the other start components summed subject to injectivity.

Use this replacement in \(K^* K\), using row \(z\) of \(K^*\) to sample \(x\). For the set \(S\) of positions occupied in \(x\), we have for every fixed set \(F\)
\begin{equation}
 \Pr(S\supset F)\le (k/N)^{|F|}.                                                            \label{ten:weighted:eq:2}
\end{equation}
In fact the upper product-of-marginals bounds for inclusions hold at the start of \(K^*\), which uses a reverse scan, from a fixed tuple \(z\). Such bounds are preserved by a random switch: for an inclusion testing just one of the switched sites, average the two bounds, and for one testing both, their product of marginals at these sites can only increase under averaging. At the end, the marginals are \(k/N\), since a single position is exactly uniform after all bit updates. This is the particular negative-dependence property needed here; stronger preservation results for partial symmetrizations are proved in \cite[Theorems~4.9 and~4.20]{BorceaBrandenLiggett2009}.

For a given \(B\), to bound the row sum using the term with weight product on \(B\) and the no-isolated-vertex condition, use true sweep paths for \(B\) and independent paths with uniform ends for the other components (we can disregard the final injectivity restriction for upper bounding the probability). If \(w\ne 1\), the two specified paths go through the same switch at a stage. We may pre-sample, independently of \(S\), two options at every starting position: a path using a common sweep permutation with its specified routing, and a path from independent uniform bits (independent paths for this option at different positions). Thus we only need upper bound the probability of a no-isolated-vertex interaction graph on \(S\) using some assignment of these types.

Such a graph contains a spanning forest with \(c\le k/2\) rooted trees of size at least 2, which we may take as ordered trees arranged in order of their root positions. The number of forest shapes needed for each \(c\) is at most \(C_0^k\), for an absolute \(C_0\) (e.g. by the traversal encoding of plane forests). Choose distinct roots in at most \(\binom Nc\) ways and path types in at most \(2^k\) ways. For these choices the expected number of embeddings with distinct positions and meetings along the tree edges is at most \((2d)^{k-c}\). To see this, condition on the full common routing and enumerate children after their parents, exposing independent paths as needed. For a child of common routing type, at each possible meeting stage there are at most two starting positions since the full routing sends exactly two paths through the switch. For a child of independent type, summed over fresh positions the expected number meeting a given parent at stage \(i\) is at most \(2^i 2^{-(i-1)}\): the suffix after \(i\) at the start must agree with the parent path there, and the new prefix of length \(i-1\) must agree. These estimates apply to the sums of extension probabilities conditional on the paths for a partial embedding, by independence for unused positions of the independent type. Each embedding uses \(k\) sites, so we may multiply the expected counts by \((k/N)^k\) by \eqref{ten:weighted:eq:2}. Summing over \(B\), shapes and embeddings gives, for the signed kernel \(K^*D\), an absolute row sum bounded by
\[
 (C_1 d)^k (k/N)^{k/2},
\]
using \(\binom Nc\le (C_2 N/k)^{k/2}\) for \(c\le k/2\). An empty set of such forests gives zero probability. The column sum bound for \(D\) still works for \(K^*D\). The row-column sum bound on the \(\ell^2\) operator norm now bounds \(K^*K\) on the specified subspace, proving \eqref{ten:weighted:eq:1}.


\subsection{The bounded-exponent induction}
\label{ten:weighted:part:271}
\begin{proof}[Proof of Theorem~\ref{thm:weighted}]
All sweep matrices are contractions. Partitions of length greater than \(N/2\) need no estimate due to the matching projection, and the trivial one satisfies \eqref{eq:weighted-main} at any exponent. For \(1\le k=N-\lambda_1\le N^{1-\epsilon}\), \eqref{ten:weighted:eq:1} and \(W_\lambda\le D_\lambda\le N^k\) suffice at any exponent at least an absolute sufficiently large constant, for sufficiently large \(d\). Indeed the norm bound \eqref{ten:weighted:eq:1} in this range is at most \(N^{-\epsilon k/16}\) for large \(d\).

For the remaining partitions, suppose \eqref{eq:weighted-main} holds at the two smaller bit sizes in the split. Use \eqref{ten:weighted:eq:20} at \(p=\max(p_{d_1},p_{d_2})\), with the subgroup inputs given by the two subgroup sweep parts. Each input on a tuple of partitions is a tensor product of smaller sweep matrices. Hence each weighted \(p\)-th power term of the inputs is at most 1, and the sums again use only the tuple count bounds. We get
\[
 W_\lambda \|K_d(\lambda)\|_p^p \le \exp(C (\log D_\lambda)/d)
\]
with \(C\) absolute, for all sufficiently large \(d\). In particular \(\|K_d(\lambda)\|_\infty^p\le D_\lambda^{-1/2}\) for these \(d\), by the lower bound on \(W_\lambda\). Increasing \(p\) to \(p(1+C'/d)\) for a large enough absolute \(C'\) gives \eqref{eq:weighted-main}, by using this operator norm on the additional powers.

To implement the induction, take an absolute base threshold on \(d\) sufficiently high for all these estimates. For each of the finitely many base dimensions a nontrivial sweep matrix has operator norm strictly less than 1: it is a product of orthogonal projections from the matchings, and equality of norms on a unit vector would require a common invariant vector for the matching switch subgroups. These generate the symmetric group since all individual edge transpositions of a connected cube are included. Thus a common sufficiently large base exponent can be used for \eqref{eq:weighted-main}, taken large enough for the sparse range as well. Thereafter use the recurrence just given with factor \(1+C'/d\) and the maximum of the two exponents for the split. These exponents are bounded above by an absolute constant since the dimensions in bits along a descent are successively halved up to rounding, so the products of the increment factors stay bounded. This proves the theorem.\end{proof}


\subsection{Full-deck mixing}
\label{ten:weighted:part:285}
In particular for nontrivial \(\lambda\) with length at most \(N/2\) we have
\[
 \|K_d(\lambda)\|_\infty\le D_\lambda^{-(1-b_0)/p_*}.
\]
For an absolute sufficiently large number \(j\) of sweeps, by Plancherel and Cauchy--Schwarz \cite[Sections~2 and~3]{DiaconisShahshahani1981}, the squared total variation distance to uniform of their product law is at most
\[
 \frac14\sum_{\lambda\ne(N)} D_\lambda \|K_d(\lambda)^j\|_2^2
 \ \le\ \frac14\sum_{\substack{\lambda\ne(N)\\ \lambda^{\rm tr}_1\le N/2}} D_\lambda^{-2},
\]
where we took \(j\) large enough for the last exponent, using the operator norm to bound the Hilbert-Schmidt norm. The sum bound tends to zero. Indeed for \(k>N^{1-\epsilon}\) this follows from \eqref{ten:weighted:eq:3} and the partition count, and for \(1\le k\le N^{1-\epsilon}\), \(D_\lambda\ge e^{-1}(N/k)^k\) for large \(N\). The latter follows by the top-row hook estimate in the proof of Lemma~\ref{lem:surviving-dimension}, now with \(k\) boxes outside the row. There are at most \(2^k\) partitions of the tail of size \(k\). Since the permutation product law determines the deck from any fixed starting deck, the bound is a worst-case bound. Thus a fixed number $j$ of sweeps gives total variation distance at most $1/4$ for all sufficiently large $d$. For each of the finitely many remaining $d$, the strict nontrivial Fourier norm gap established above gives convergence under repeated forward sweeps. Increasing $j$ to cover these sizes yields one absolute sweep count for every $d\ge1$; further convolution cannot increase total variation distance. As a sweep consists of $d$ physical shuffles, this proves the upper bound in Corollary~\ref{cor:optimal-mixing}.

\begin{proposition}[Support obstruction]\label{ten:support}
Let $q_d$ be the law of one physical Thorp shuffle. For every integer $t\ge0$, its $t$-fold convolution satisfies
\[
 \|q_d^{*t}-U_{S_N}\|_{\rm TV}\ge1-\frac{2^{tN/2}}{N!}.
\]
Thus mixing to distance $1/4$ requires at least
$\lceil(2/N)\log_2(3N!/4)\rceil=2d-O(1)$ physical shuffles.
\end{proposition}
\begin{proof}
The $tN/2$ fair bits allow at most $2^{tN/2}$ permutations. Testing the support against uniform measure proves the inequality. The threshold follows by rearrangement, and Stirling's formula gives its asymptotic form.
\end{proof}
